\section{Exponential coordinates, coframes, and metrics to all orders}\label{sec:geometry}
The differential identity has an intrinsic geometric meaning. This section also isolates the corrections needed in the geometric paragraph of the source article.

\subsection{The Maurer--Cartan form and its complete coordinate expansion}
Let $G$ be a finite-dimensional Lie group, with basis $e_1,\ldots,e_d$ of $\gfrak$ and structure constants
\begin{equation}
 [e_i,e_j]=f_{ij}{}^k e_k.
\end{equation}
Repeated upper--lower indices are summed. The left Maurer--Cartan form is the $\gfrak$-valued one-form $\theta_g=(L_{g^{-1}})_*$; in a matrix group it is $g^{-1}\dd g$. For $g=\exp X$, where $X=X^i e_i$, \cref{eq:left-dexp} gives
\begin{align}
 \theta&=\sum_{n=0}^{\infty}\frac{(-1)^n}{(n+1)!}\ad_X^n(\dd X)
       =e_a W^a{}_j(X)\,\dd X^j,\label{eq:theta-all}\\
 M^a{}_j&=X^i f_{ij}{}^a,\qquad
 W=f(M)=\sum_{n=0}^{\infty}\frac{(-1)^nM^n}{(n+1)!}.
 \label{eq:W-all}
\end{align}
The first terms are, with all contractions visible,
\begin{equation}\label{eq:theta-first}
 \theta=\dd X^i e_i
 -\tfrac12 X^i\dd X^j f_{ij}{}^k e_k
 +\tfrac16 X^iX^j\dd X^k f_{jk}{}^\ell f_{i\ell}{}^m e_m
 -\cdots.
\end{equation}
For arbitrary $n\geq1$ the contraction in the $n$th term is specified by
\begin{equation}
 (M^n)^a{}_j=
 X^{i_1}\cdots X^{i_n}
 f_{i_nj}{}^{b_1}f_{i_{n-1}b_1}{}^{b_2}\cdots
 f_{i_1b_{n-1}}{}^a,
\end{equation}
with the evident interpretation when $n=1$. Thus \cref{eq:theta-all} is an all-terms formula, not merely an instruction to continue a pattern.

The notation
\begin{equation}\label{eq:W-quotient}
 W=\frac{I-e^{-M}}{M}
\end{equation}
means the entire function $f(M)$, not multiplication by an existing inverse of $M$. In fact
\begin{equation}\label{eq:M-singular}
 M X=\ad_X X=0,\qquad M(0)=0,\qquad W(0)=I.
\end{equation}
In positive dimension $M$ is therefore always singular. The matrix supplying the coframe in exponential coordinates is \emph{$W$, not $M$}. Near the identity its dual frame has coefficient matrix $W^{-1}=h(M)$. More generally,
\begin{equation}\label{eq:W-invertibility}
 W\text{ is invertible}\quad\Longleftrightarrow\quad
 \spec(M)\cap\{2\pi\ii k:k\in\mathbb Z\setminus\{0\}\}=\varnothing.
\end{equation}
Indeed, complex triangularization makes the eigenvalues of $f(M)$ equal to $f(\lambda)$, counted with algebraic multiplicity. The zeros of $f$ are precisely the indicated nonzero points. This argument also explains why the zero eigenvalue forced by \cref{eq:M-singular} is harmless for $W$.

\begin{proposition}[Maurer--Cartan equation]\label{prop:maurer-cartan}
The form $\theta$ satisfies
\begin{equation}\label{eq:maurer-cartan}
 \dd\theta+\tfrac12[\theta,\theta]=0.
\end{equation}
Here $\tfrac12[\theta,\theta](U,V)=[\theta(U),\theta(V)]$.
\end{proposition}
\begin{proof}
On left-invariant vector fields $U,V$, both $\theta(U)$ and $\theta(V)$ are constant, so the definition of exterior differentiation gives
$\dd\theta(U,V)=-\theta([U,V])=-[\theta(U),\theta(V)]$.
Left-invariant vectors span each tangent space, proving the identity. In a matrix realization the same proof is the calculation
$\dd(g^{-1})=-g^{-1}(\dd g)g^{-1}$ and
$\dd(g^{-1}\dd g)=-(g^{-1}\dd g)\wedge(g^{-1}\dd g)$.
\end{proof}
In particular, for a two-parameter map $g(s,t)$, writing
$A=\theta(\partial_sg)$ and $U=\theta(\partial_tg)$ gives
$\partial_sU-\partial_tA+[A,U]=0$. This is the mixed-variation identity used in \cref{sec:local-groups}; it involves no matrix embedding assumption.

\subsection{Pullback tensors: the necessary Jacobian and nondegeneracy condition}
Let $B$ be a symmetric bilinear form on $\gfrak$. Left translation defines a symmetric tensor on $G$ by
$\langle u,v\rangle_g=B(\theta_g u,\theta_g v)$.
It is a pseudo-Riemannian metric exactly when $B$ is nondegenerate. Given a smooth map $\varphi:N\to G$, use coordinates $q^\alpha$ and write locally
$\varphi(q)=\exp(X(q))$. Set
\begin{equation}\label{eq:pullback-frame}
 E^a{}_\alpha=W^a{}_j(X(q))\,\partial_\alpha X^j.
\end{equation}
The pullback tensor is
\begin{equation}\label{eq:pullback-metric}
 g_{\alpha\beta}(q)=B_{ab}E^a{}_\alpha E^b{}_\beta.
\end{equation}
This is immediate by applying the definition of the pullback to the coordinate vectors. For an arbitrary, rather than left-invariant, metric on $G$, its coefficients $B_{ab}$ in the left-invariant frame depend on $\varphi(q)$; the same formula holds with this dependence retained. Only when the coordinates themselves are $q^i=X^i$ does the Jacobian disappear and the expression reduce to $g_{ij}=B_{ab}W^a{}_iW^b{}_j$.

Three distinct qualifications are important. An arbitrary pullback need not be nondegenerate: a constant map gives the zero tensor. If $B$ is positive definite, an immersion has a positive-definite pullback; for indefinite $B$, even an immersion can have a degenerate pullback. For example, the immersion $t\mapsto(t,t)$ into $(\RR^2,\dd x^2-\dd y^2)$ pulls the metric back to zero. Finally, the Killing form
\begin{equation}\label{eq:killing}
 B_K(U,V)=\tr(\ad_U\ad_V)
\end{equation}
need not be nondegenerate. In an abelian Lie algebra it vanishes identically. Consequently, choosing it does not automatically give a metric or a vielbein construction.

For completeness, the Killing form is symmetric by cyclicity of trace, and it is infinitesimally invariant:
\begin{equation}\label{eq:B-invariant}
 B_K([X,U],V)+B_K(U,[X,V])=0.
\end{equation}
To verify this, use $\ad_{[X,U]}=[\ad_X,\ad_U]$ and combine the four trace terms into the trace of a commutator. The latter identity for adjoint maps is the Jacobi identity. This proves the invariance used next without assuming nondegeneracy.

\subsection{All-order metric and volume formulas}
Put $U_\alpha=\partial_\alpha X$. Substituting \cref{eq:W-all} in \cref{eq:pullback-metric} gives the complete expansion
\begin{equation}\label{eq:metric-double}
 g_{\alpha\beta}
 =\sum_{r,s\geq0}\frac{(-1)^{r+s}}{(r+1)!(s+1)!}
 B(M^rU_\alpha,M^sU_\beta).
\end{equation}
In finite dimension this double series converges absolutely for every $M$ and all vectors, regardless of whether the tensor it defines is nondegenerate.

\begin{proposition}[Invariant-form simplification]\label{prop:metric-even}
Suppose $B([X,U],V)=-B(U,[X,V])$. Then
\begin{align}
 g_{\alpha\beta}
 &=B\left(U_\alpha,\frac{2(\cosh M-I)}{M^2}\,U_\beta\right)\label{eq:metric-entire}\\
 &=\sum_{m=0}^{\infty}\frac{2}{(2m+2)!}
 B(U_\alpha,M^{2m}U_\beta).\label{eq:metric-all}
\end{align}
The quotient denotes its entire extension. The first coefficients are $1$, $1/12$, $1/360$, $1/20160$, and so on.
\end{proposition}
\begin{proof}
Move each $M$ in the first argument across $B$, picking up a minus sign. Thus
$B(f(M)U,f(M)V)=B(U,f(-M)f(M)V)$.
The scalar identity
$f(-z)f(z)=(e^z+e^{-z}-2)/z^2$
then proves both formulas by power-series multiplication.
\end{proof}
Notice that all odd orders disappear only under the stated invariance condition; no such cancellation was assumed in \cref{eq:metric-double}.

When $B$ is nondegenerate and $X^1,\ldots,X^d$ are valid exponential coordinates, the metric matrix is $W^TBW$. Its volume density is consequently
\begin{equation}\label{eq:volume}
 \sqrt{\abs{\det g}}=\abs{\det W}\sqrt{\abs{\det B}},\qquad
 \det W=\prod_{\lambda\in\spec M}\frac{1-e^{-\lambda}}{\lambda},
\end{equation}
where the factor at zero is $1$ and multiplicities are included. Near $M=0$ its logarithm has a second complete expansion:
\begin{equation}\label{eq:log-volume}
 \log\det W=-\tfrac12\tr M+
 \sum_{m=1}^{\infty}\frac{b^+_{2m}}{2m}\tr(M^{2m}).
\end{equation}
For instance the first even terms are $(\tr M^2)/24-(\tr M^4)/2880$.
To prove this, differentiate the scalar logarithm:
\begin{equation}
 \frac{\dd}{\dd z}\log f(z)
 =-\tfrac12+\sum_{m\geq1}b^+_{2m}z^{2m-1}.
\end{equation}
Integration and $f(0)=1$ give its Taylor series; triangularization and summing over eigenvalues yield \cref{eq:log-volume}. For example, $\norm M<2\pi$ is a sufficient condition for this Taylor series. The determinant formula \cref{eq:volume} itself has no Taylor-radius restriction.
